\documentclass[aps,prl,reprint]{revtex4-2}

\usepackage{graphicx}
\usepackage{amsmath,amssymb}

\begin{document}
% Paste after the existing Appendix A and before the bibliography.
% The manuscript already enters appendix mode; do not add another \appendix.
% Requires only amsmath, already used in the manuscript.

\section{Surrogate data generation, training, and testing}
\label{app:surrogate_training}

This appendix describes the dataset and fitting procedures for the neural
network (NN) predicting the electron cutting-velocity profile and the support
vector machine (SVM) classifying GYRAZE convergence. The physical model and
the coupling to Gkeyll are described in the main text.

\subsection{Sampling and data preprocessing}

A total of 3318 GYRAZE runs were attempted with uniformly sampled inputs
$\mathbf{x}=(\alpha,\gamma,\Phi)$, where
$\gamma=\rho_e/\lambda_d$ and
$\Phi=e(\phi_{\mathrm{mpe}}-\phi_w)/T_e$. The sampling domain was
\begin{equation}
  \alpha\in[0^\circ,10^\circ],\qquad
  \gamma\in[0,9],\qquad
  \Phi\in[0,10].
\end{equation}
The runs were distributed across 16 CPU cores. Of these runs, 1566 converged
and supplied input/cutting-velocity profile pairs $(\mathbf{x}, v_{\parallel \text{cut}}(\mu_{\text{grid}}))\in \mathbb{R}^3 \times \mathbb{R}^{20}$ for NN regression. This provided the following approximate range of NN training data, with sampling distribution determined by the convergence behavior of GYRAZE
\begin{equation}
  \alpha\in[2^\circ,10^\circ],\qquad
  \gamma\in[0.5,4],\qquad
  \Phi\in[1,10].
\end{equation}

All 3318 runs
supplied convergence labels ('yes','no') for SVM classification. These labels describe
the numerical convergence of GYRAZE under the specified solver settings;
nonconvergence does not establish the absence of a physical sheath solution. Native GYRAZE convergence tolerances were used: $<0.1\%$ MP mean residual, $<1\%$ DS mean residual, and $<0.5\%$ relative current error. A maximum of 1000 iterations was used for each phase, with a timeout of 3 minutes per sample.

For each learning task, the relevant dataset was randomly partitioned into
training, validation, and test subsets in an approximate 60:20:20 ratio. The input data was normalized by the mean and sample variance of the training dataset.
Training subsets were used for parameter fitting, validation subsets for
hyperparameter selection, and test subsets for final performance evaluation.
\subsection{Neural-network architecture and training}

The NN maps the three input parameters to the dimensionless
cutting-velocity profile $v_{\mathrm{NN}}(\mathbf{x};\mu)$ used in the main
text. Only converged GYRAZE runs were used to fit and evaluate this mapping.
The selected model was a fully connected feedforward network with SiLU
hidden activations and a linear output layer. The input dimension is 3, there are 3 hidden layers in the network of width 75, and the output dimension is 20 (corresponding to the native discretization $\mu_\text{grid}$ in GYRAZE):
\[
3 \rightarrow{75}\rightarrow{75}\rightarrow{75}\rightarrow{20}
\]

Network parameters were fitted using the Adam optimizer with a
mean-squared-error loss on mini-batches. Activation function, depth, and
hidden-layer width were varied coarsely, and the configuration with the lowest
validation loss was selected. A learning rate $\eta = 1\times 10^{-5}$, batch size $32$, and 500 epochs were used in training.

On the held-out converged runs, the selected NN achieved
$R^2=0.9991$ and a test MSE of $8.34\times10^{-4}$.

\begin{figure}[h]
\centering
\includegraphics[width=\columnwidth]{ConvNNprof.png}
\caption{Predictive performance of the Convergent NN Surrogate on the testing dataset.}
\label{fig:conv_nn_prof}
\end{figure}

\subsection{SVM convergence classification}

A soft-margin SVM with a radial-basis-function kernel was trained on the convergence
labels of all attempted runs, using the training/validation/test procedure
described above. The selected hyperparameters were $C=0.1$ and kernel
coefficient $\kappa=3.16$, where
\begin{equation}
 K(\mathbf{z},\mathbf{z}')=
 \exp\!\left[-\kappa\|\mathbf{z}-\mathbf{z}'\|_2^2\right],
\end{equation}
and $\mathbf{z}$ denotes the input representation supplied to the
classifier. Hyperparameters were selected
by validation accuracy. The selected classifier achieved a held-out test
accuracy of $96.62\%$.
\begin{figure*}[h]
\centering
\includegraphics[width=16cm]{svm_boundary_split.png}
\caption{Predictive performance of the SVM on the full dataset. Precision is high compared to recall, which suits our application; the convergence region should admit few nonconverged runs, characterizing parameters with consistent convergence in GYRAZE. Points inside the decision boundary (left) and points outside the decision boundary (right).}
\label{fig:conv_nn_prof}
\end{figure*}

\subsection{Use outside the GYRAZE convergence region}

[to be discussed]


\end{document}